\usepackage[spanish,es-tabla]{babel}
\usepackage[top=2cm,bottom=2cm,right=2cm,left=2cm]{geometry}
\usepackage{amsmath}
\usepackage{amssymb}
\usepackage{amsthm}
\usepackage{empheq}
\usepackage{mdframed}
\usepackage{booktabs}
\usepackage{lipsum}
\usepackage{graphicx}
\usepackage{color}
\usepackage{psfrag}
\usepackage{pgfplots}
\usepackage{bm}
\usepackage{multicol}
\usepackage{wrapfig}
\usepackage{float}
\usepackage{wrapfig}
\usepackage{multirow}
\usepackage{biblatex}
\usepackage{csquotes}
\usepackage{parskip}
\usepackage{cancel}
\usepackage{listing}
\usepackage{tikz}
\usepackage{verbatim}
\usepackage{subcaption}
\usepackage{enumitem}
\usepackage{xcolor}
\usepackage{xfrac}
\usepackage[load-configurations = abbreviations]{siunitx}
\usepackage{fix-cm}
\usepackage{yhmath}
\usepackage[bookmarks=true,colorlinks=true,linkcolor=black,citecolor=black,menucolor=black,urlcolor=black,hidelinks,breaklinks=true]{hyperref}
\usepackage{mathrsfs}
\usepackage{tcolorbox}
\usepackage{amsfonts}
\usepackage[bb=ams]{mathalpha}
\usepackage{mathrsfs}
\usepackage{xstring}

\graphicspath{{fig/}}
\addbibresource{biblio.bib}
\usepgfplotslibrary{colorbrewer}
\pgfplotsset{width=8cm,compat=1.9}
\definecolor{negro}{RGB}{0,0,0}
\definecolor{azul}{RGB}{0,0,255}
\definecolor{verde}{RGB}{0,255,0}
\definecolor{cian}{RGB}{0,255,255}
\definecolor{rojo}{RGB}{255,0,0}
\definecolor{magenta}{RGB}{255,0,255}
\definecolor{amarillo}{RGB}{255,255,0}
\definecolor{blanco}{RGB}{255,255,255}

\definecolor{limegreen}{cmyk}{0.75,0,1,0}
\definecolor{orange}{cmyk}{0,0.8,0.95,0}
\definecolor{purple}{cmyk}{0.75,1,0,0}

\newcommand{\uvec}[1]{~ \boldsymbol{\hat{\textbf{#1}}}}

\newsavebox\foobox
\newlength{\foodim}
\newcommand{\slantbox}[2][0]{\mbox{%
        \sbox{\foobox}{#2}%
        \foodim=#1\wd\foobox
        \hskip \wd\foobox
        \hskip -0.5\foodim
        \pdfsave
        \pdfsetmatrix{1 0 #1 1}%
        \llap{\usebox{\foobox}}%
        \pdfrestore
        \hskip 0.5\foodim
}}
\def\Laplacesymb{\slantbox[-.45]{$\mathscr{L}$}}
\def\inverseLaplacesymb{\Laplacesymb^{-1}}
\def\Fouriersymb{\slantbox[-.45]{$\mathscr{F}$}}
\def\Fouriersinsymb{\Fouriersymb_{\textup{s}}}
\def\Fouriercossymb{\Fouriersymb_{\textup{c}}}
\def\inverseFouriersymb{\Fouriersymb^{-1}}

\newcommand{\Laplace}[2]{
    \IfEqCase{#1}{
        {}{\Laplacesymb\left\lbrace#2\right\rbrace}
        {i}{\inverseLaplacesymb\left\lbrace#2\right\rbrace}
    }[\PackageError{tree}{Undefined option to tree: #1}{}]
}
\newcommand{\Fourier}[2]{
    \IfEqCase{#1}{
        {}{\Fouriersymb\left\lbrace#2\right\rbrace}
        {i}{\inverseFouriersymb\left\lbrace#2\right\rbrace}
        {sin}{\Fouriersinsymb\left\lbrace#2\right\rbrace}
        {cos}{\Fouriercossymb\left\lbrace#2\right\rbrace}
    }[\PackageError{tree}{Undefined option to tree: #1}{}]
}

\sisetup{
  detect-family,
  per-mode = symbol,
  separate-uncertainty = true,
  inter-unit-product = \cdot,
  %output-decimal-marker = {,},
  %exponent-product = \cdot
}

\begin{document}

\onecolumn\begin{@twocolumntrue}
    \begin{minipage}{0.3\textwidth}{\includegraphics[scale=0.24]{Escudo_UD.pdf}} % Cambiar escudo
    \end{minipage}
    \vspace{10pt}
    \begin{minipage}{0.677\textwidth}
        \begin{center}
            \vspace{12mm}
        
            \Large{\textbf{<<Título>>}}
            \vspace{3mm}
            
            \large{\textbf{<<Autores>>}}
            \vspace{2mm}
    
            \large{\textit{<<Carrera>>, <<Universidad>>}}
            \vspace{1mm}
            
            <<Fecha>> % Formato (mes) de (año)
        \end{center}
\vspace{5pt}
\end{minipage}

\centerline{\rule{0.95\textwidth}{0.4pt}} % Línea negra
\end{@twocolumntrue}

\begin{multicols}{2}

\section{Periodo de una función}

\begin{enumerate}[leftmargin=15pt]
    \item Encontrar el periodo de las siguientes funciones:
    \begin{enumerate}
        \item $f(t) = cos \left(n \pi t\right)$ 
        \item $f(t) = \sin \left(t\right) + \sin \left(\frac{t}{3}\right) + \sin \left(\frac{t}{5}\right)$
    \end{enumerate} 
\end{enumerate}

\subsubsection*{\textbf{Solución:}}


Para la solución de este ejercicio, se tiene en cuenta la forma general para hallar el periodo de una función trigonométrica, por tanto, para toda función trigonométrica con periodo natural $T_n$, se tiene su periodo $T_f$ de la siguiente manera:
\begin{equation}
    T_f = \frac{T_n}{B}
\label{eq:1}
\end{equation}
Dónde $B$ yace de la forma general de una función trigonométrica, la cuál es:
\begin{equation*}
    f(t) = A\sin\left(Bx + C\right) + D
\end{equation*}
Forma usada para cualquier función trigonométrica no necesariamente para la función seno, de esta manera, solo analizando la forma de la función y además conociendo el periodo natural de la función que estemos utilizando, podremos conocer el periodo de dicha función.

Los periodos naturales de las diferentes funciones trigonométricas son:

\begin{itemize}
    \item Para \textit{seno, coseno, secante} y \textit{cosecante} es $2\pi$.
    \item Para \textit{tangente} y \textit{cotangente} es $\pi$.
\end{itemize}

Para una superposición de funciones, se debe buscar el mínimo común múltiplo (m.c.m.) de esta manera se busca en qué instante todas las funciones pertenecientes a la superposición están en fase, o lo que sería igual a que comparten el mismo periodo.

Teniendo todo esto en cuenta y aplicándolo, resulta:

$a) ~ f(t) = \cos \left(n \pi t\right)$

Tomando (\ref{eq:1}) y teniendo $B = n \pi$ y $T_n = 2\pi$, el periodo $T_f$ de la función es:
\begin{equation}
    T_f = \frac{2 \cancel{\pi}}{n \cancel{\pi}} \Rightarrow \textcolor{rojo}{T_f = \frac{2}{n}}
    \label{res:1}
\end{equation}
Obteniendo así que el periodo de la función es \textcolor{rojo}{$T_f = \sfrac{2}{n}$}.

$b) ~ f(t) = \sin \left(t\right) + \sin \left(\frac{t}{3}\right) + \sin \left(\frac{t}{5}\right)$

Tomando (\ref{eq:1}) nuevamente y teniendo $B_1 = 1$, $B_2 = \left(\frac{1}{3}\right)$, $B_3 = \left(\frac{1}{5}\right)$ y $T_n = 2\pi$ para todas las funciones superpuestas, definimos $T_f$ para cada una y posteriormente hallamos el m.c.m. entre ellas, así de esta manera obtener el periodo general de la función, por tanto:
\begin{align}
    T_{f1} & = \frac{2\pi}{1} \Rightarrow T_{f1} = 2\pi \\
    T_{f2} & = \frac{2\pi}{ \frac{1}{3} } \Rightarrow T_{f2} = 6\pi \\
    T_{f3} & = \frac{2\pi}{ \frac{1}{5} } \Rightarrow T_{f3} = 10\pi
\end{align}
Usando los resultados obtenidos en (3), (4) y (5)
\begin{center}

$\begin{array}{c c c | c}
    2\pi & 6\pi & 10\pi & 2\\
    1\pi & 3\pi & 5\pi & 3\\
    \vdots & 1\pi & 5\pi & 5\\
    \vdots & \vdots & 1\pi & \pi\\
    1 & 1 & 1 & \hphantom{1}\\ \hline
    \hphantom{1\pi} & \hphantom{1\pi} & \hphantom{1\pi} & 30\pi
\end{array}$

\end{center}
De esta manera se concluye qué:
\begin{equation}
    \textcolor{rojo}{T_f = 30\pi}
\label{res:2}
\end{equation}

\begin{enumerate}[leftmargin=15pt, resume]
    \item Si $f(t)$ es una función  periódica de $t$ con periodo $T$, demostrar que $f(at)$, para $a \neq 0 $, es una función periódica de $t$ con periodo $\frac{T}{a}$.
\end{enumerate}

\subsubsection*{\textbf{Solución:}}

A la hora de demostrar que una función $f(t)$ posee un periodo $T$ y lo que implicaría que tiene la forma $f(t+T)$, decimos que para demostrar que una función cualquiera sea periódica esta debe tener la misma forma mostrada previamente, así, al variar los parámetros correspondientes ya sea insertando corregimientos o eliminando, el resultado final debe ser el parámetro más su periodo $T$, por tanto, para una función $f(at)$, se dice que tiene un periodo $\sfrac{T}{a}$ de tal manera que al operar

\begin{equation*}
    f(at) = f\left(a\left(t + \frac{T}{a}\right)\right) = f\left(at + T\right)
\end{equation*}

Donde si se toma $T_f = \sfrac{T}{a}$, ya que anteriormente se dice que ese es el periodo de la función, entonces, queda demostrado que para una función $f(at)$ su periodo es $\sfrac{T}{a}$, resultando en el periodo normal de la función $T_f \equiv T$.

\section{Series de Fourier}

\begin{enumerate}[leftmargin=15pt]
    \item Demostrar que la serie de Fourier para la función $F(t)$ definida por:
\end{enumerate}

\begin{equation}
    F(t) =
    \begin{cases}
    0,&\text{si } \sfrac{-\pi}{2} \leq t < 0\\
    1,&\text{si } \hspace{18pt} 0 \leq t < \sfrac{\pi}{2}
    \end{cases}
\label{enu:1}
\end{equation}
\hspace{12pt} y
\begin{equation*}
    F(t) = F\left(t + \pi\right)
\end{equation*}

\begin{figure}[H]
\begin{center}
\includegraphics[scale=0.27]{fig0.png}
\caption{Gráfica de $F(t)$ de la forma (\ref{enu:1}) con periodo $T_f = \pi$}
\label{figenu:1}
\end{center}
\end{figure}
\hspace{12pt} Viene dada por:

\begin{equation}
    F(t) = \frac{1}{2} + \frac{2}{\pi} \sum_{1}^{\infty} \frac{\sin \left(2\left(2n-1\right)t\right)}{\left(2n-1\right)}
    \label{enu:2}
\end{equation}

\subsubsection*{\textbf{Solución:}}

Para deducir la serie de Fourier trigonométrica de cualquier función continúa, suave de la forma $F\left(t+T\right)$ se plantea la forma general para la serie, la cual es:
\begin{align*}
    F(t) = \frac{a_0}{2} + &\sum_{n = 1}^{\infty} a_n \cos\left(n \omega_f t\right) + b_n \sin\left(n \omega_f t\right)\\
    &~\mathrm{donde}~ \omega_f = \frac{2\pi}{T_f}
\end{align*}
De la cual se deducen los valores de los coeficientes $a_0,~ a_n~ \mathrm{y}~ b_n$; esto gracias a una serie de cálculos, por tanto
\begin{equation}
    a_0 = \frac{2}{T_f} \int_{\sfrac{-T_f}{2}}^{\sfrac{T_f}{2}} F(t) \,dt
    \label{eq:2}
\end{equation}
\begin{equation}
    a_n = \frac{2}{T_f} \int_{\sfrac{-T_f}{2}}^{\sfrac{T_f}{2}} F(t) \cos\left(n \omega_f t\right) \,dt
    \label{eq:3}
\end{equation}
\begin{equation}
    b_n = \frac{2}{T_f} \int_{\sfrac{-T_f}{2}}^{\sfrac{T_f}{2}} F(t) \sin\left(n \omega_f t\right) \,dt
    \label{eq:4}
\end{equation}
De esta manera, reemplazando en (\ref{eq:2}), (\ref{eq:3}) y (\ref{eq:4}) para hallar los respectivos coeficientes, teniendo $T_f = \pi$ resulta:

Para $a_0$:
\begin{align*}
    a_0 &= \frac{2}{\pi} \int_{\sfrac{-\pi}{2}}^{\sfrac{\pi}{2}} F(t) \,dt\\
    &= \frac{2}{\pi} \left[\int_{\sfrac{-\pi}{2}}^{0} 0 \,dt + \int_{0}^{\sfrac{\pi}{2}} 1 \,dt\right]\\
    &= \frac{2}{\pi} \left(\int_{0}^{\sfrac{\pi}{2}}\,dt\right)\\
    &= \frac{2}{\pi} \left(t |^{\sfrac{\pi}{2}}_{0}\right)\\
    &= \frac{2}{\pi} \left(\frac{\pi}{2} - 0\right)\\
    &= \frac{\cancel{2}}{\cancel{\pi}} \left(\frac{\cancel{\pi}}{\cancel{2}}\right)\\
    &= 1
\end{align*}
Resultando
\begin{equation}
    a_0 = 1
    \label{res:3}
\end{equation}
Para $a_n$, teniendo $\omega_f = 2$:
\begin{align*}
    a_n &= \frac{2}{\pi} \int_{\sfrac{-\pi}{2}}^{\sfrac{\pi}{2}} F(t) \cos\left(2nt\right) \,dt\\
    &= \frac{2}{\pi} \left[\int_{\sfrac{-\pi}{2}}^{0} 0 \cos\left(2nt\right) \,dt + \int_{0}^{\sfrac{\pi}{2}} 1 \cos\left(2nt\right) \,dt\right]\\
    &= \frac{2}{\pi} \left[\int_{0}^{\sfrac{\pi}{2}}\cos\left(2nt\right) \,dt\right]\\
    &= \frac{2}{\pi} \left[\frac{1}{2n} \left(\sin\left(2nt\right) |^{\sfrac{\pi}{2}}_{0}\right)\right]\\
    &= \frac{2}{\pi} \left[\frac{1}{2n} \left(\sin\left(\cancel{2}n\frac{\pi}{\cancel{2}}\right) - \sin\left(2n(0)\right)\right)\right]\\
    &\mathrm{Ya~que}~\sin\left(n\pi\right) = \sin\left(0\right) = 0\\
    &= \frac{2}{\pi} \left[0\right]\\
    &= 0
\end{align*}
Resultando
\begin{equation}
    a_n = 0
    \label{res:4}
\end{equation}
Para $b_n$, teniendo $\omega_f = 2$:
\begin{align*}
    b_n &= \frac{2}{\pi} \int_{\sfrac{-\pi}{2}}^{\sfrac{\pi}{2}} F(t) \sin\left(2nt\right) \,dt\\
    &= \frac{2}{\pi} \left[\int_{\sfrac{-\pi}{2}}^{0} 0 \sin\left(2nt\right) \,dt + \int_{0}^{\sfrac{\pi}{2}} 1 \sin\left(2nt\right) \,dt\right]\\
    &= \frac{2}{\pi} \left[\int_{0}^{\sfrac{\pi}{2}}\sin\left(2nt\right) \,dt\right]\\
    &= \frac{2}{\pi} \left[-\frac{1}{2n} \left(\cos\left(2nt\right) |^{\sfrac{\pi}{2}}_{0}\right)\right]\\
    &= \frac{2}{\pi} \left[-\frac{1}{2n} \left(\cos\left(\cancel{2}n\frac{\pi}{\cancel{2}}\right) - \cos\left(2n(0)\right)\right)\right]\\
    &\mathrm{Ya~que}~\cos\left(n\pi\right) = \left(-1\right)^{n} \mathrm{~y~} \cos\left(0\right) = 1\\
    &= \frac{2}{\pi} \left[-\frac{1}{2n}\left(\left(-1\right)^{n} - 1\right)\right]
\end{align*}
En este caso, para una $n$ par, $\left(-1\right)^{n} = 1$ y para una impar, $\left(-1\right)^{n} = -1$. Teniendo que si $n$ es par, $b_n = 0$, lo cual es un valor inútil para nosotros, ya que daría que la serie de Fourier es igual a una constante ($\sfrac{a_0}{2}$), por ende se fuerza que $n$ sea impar, resultando $b_n$ así:
\begin{equation}
    b_n = \frac{2}{\pi} \cdot \frac{1}{\left(2n-1\right)}
    \label{res:5}
\end{equation}

Tras haber obtenido los valores de los coeficientes $a_0,~a_n$ y $b_n$, se reemplazan en la ecuación general para la serie trigonométrica se Fourier. Usando (\ref{res:3}), (\ref{res:4}) y (\ref{res:5})
\begin{equation}
    \textcolor{rojo}{F(t) = \frac{1}{2} + \frac{2}{\pi} \sum_{n = 1}^{\infty} \frac{\sin \left(2\left(2n-1\right)t\right)}{2n-1}}
    \label{res:6}
\end{equation}
Comparando (\ref{res:6}) con (\ref{enu:2}) se puede ver la similitud, quedando demostrado y con ello, solucionado el ejercicio.

\begin{enumerate}[leftmargin=15pt, resume]
    \item Utilizar la calculadora \textit{DESMOS} para calcular y graficar la serie de Fourier de las siguientes funciones:
    \begin{enumerate}
        \item $F(t) = -t$ para el intervalo entre $(-\pi,~\pi)$ y $F(t) = F\left(t+2\pi\right)$.
        \item $F(t) = e^{t}$ para el intervalo entre $(-\pi,~\pi)$ y $F(t) = F\left(t+2\pi\right)$.
        \item $F(t) = \ln |t|$ para el intervalo entre $(-\pi,~\pi)$ y $F(t) = F\left(t+2\pi\right)$.
    \end{enumerate}
\end{enumerate}

\subsubsection*{\textbf{Solución:}}

Para la solución de este ejercicio se utilizó el software \textit{DESMOS}, esto con el fin de soporte en cuánto a la visualización de la gráfica y con ello, poder calcular la serie de Fourier de manera más sencilla. Se van a colocar algunas imágenes referenciando a las gráficas de las respectivas funciones y posteriormente se va a deducir la serie de Fourier con base en ello.

$a.~ F(t) = F(t + 2\pi) = -t;~ (-\pi,~\pi)$

Usando el software \textit{DESMOS}, la gráfica para $F(t)$ junto a su serie de Fourier correspondiente evaluada hasta $n=15$ es:

\begin{figure}[H]
\begin{center}
    \includegraphics[scale=0.40]{desmos-graph1.png}
    \caption{$F(t) = -t$ con su serie de Fourier entre $(-\pi,~\pi)$}
    \label{figdes:1} 
\end{center}
\end{figure}

Como se puede ver en la Figura \ref{figdes:1}, gracias a la calculadora \textit{DESMOS}, se puede saber que $a_0 = a_n = 0$, por tanto, el único valor obtenido es $b_n$, que es valor que se va a calcular a continuación:

Para $b_n$ teniendo $\omega_f = 1$:
\begin{align*}
    b_n &= \frac{\cancel{4}}{\cancel{2}\pi} \int_{0}^{\frac{\cancel{2}\pi}{\cancel{2}}} -t\sin\left(n t\right) \,dt\\
    &= -\frac{2}{\pi} \int_{0}^{\pi} t\sin\left(n t\right) \,dt
\end{align*}
Sustituyendo al usar integración por partes
\begin{align*}
    u&=t           &  dv&=\sin\left(nt\right) dt\\
    du&=dt         &  v&=-\frac{1}{n} \cos\left(nt\right)
\end{align*}
Entonces
\begin{align*}
    b_n &= -\frac{2}{\pi} \left[-\frac{t}{n}\cos\left(nt\right) - \int \frac{1}{n} \cos\left(n t\right) \,dt\right]_{0}^{\pi}\\
    &= -\frac{2}{\pi} \left[-\frac{t}{n}\cos\left(nt\right) - \frac{1}{n} \int \cos\left(n t\right) \,dt\right]_{0}^{\pi}\\
    &= -\frac{2}{\pi} \left[-\frac{t}{n}\cos\left(nt\right) - \frac{1}{n} \left(\frac{1}{n} \sin\left(nt\right)\right)\right]_{0}^{\pi}\\
    &= -\frac{2}{\pi} \left[-\frac{t\cos\left(nt\right)}{n} - \frac{\sin\left(nt\right)}{n^{2}}\right]_{0}^{\pi}\\
    &= \frac{2}{\pi} \left[\frac{nt \cos\left(nt\right) + \sin\left(nt\right)}{n^{2}}\right]_{0}^{\pi}\\
    &= \frac{2}{\pi} \left(\frac{n\pi \cos\left(n\pi\right) + \sin\left(n\pi\right)}{n^{2}}\right) -\\
    &\hphantom{=}~ \left(\frac{n(0) \cos\left(n(0)\right) + \sin\left(n(0)\right)}{n^{2}}\right)
\end{align*}
Ya que $\sin\left(n\pi\right) = \sin\left(0\right) = 0,~ \cos\left(n\pi\right) = \left(-1\right)^{n}$ y $\cos\left(0\right) = 1$,
\begin{align*}
    &= \frac{2}{\cancel{\pi}} \left(\frac{\cancel{n}\cancel{\pi}\left(-1\right)^{n}}{n^{\cancel{2}}}\right) \hphantom{,~tomando~n~impar}\\
    &= \frac{2\left(-1\right)^{n}}{n}
\end{align*}
Resultando
\begin{equation}
    \textcolor{rojo}{b_n = 2 \frac{\left(-1\right)^{n}}{n}}
    \label{res:21}
\end{equation}
Teniendo $a_0 = a_n = 0$, y usando (\ref{res:21}), la serie de Fourier para $F(t)$ es:
\begin{equation*}
    \textcolor{rojo}{F(t) = 2 \sum_{n=1}^{\infty} \left(-1\right)^{n} \frac{\sin\left(nt\right)}{n}}
\end{equation*}

$b.~F(t) = F(t + 2\pi) = e^{t};~(-\pi,~\pi)$

Usando el software \textit{DESMOS}, la gráfica para $F(t)$ junto a su serie de Fourier correspondiente evaluada hasta $n=15$ es:

\begin{figure}[H]
\begin{center}
    \includegraphics[scale=0.40]{desmos-graph2.png}
    \caption{$F(t) = e^{t}$ con su serie de Fourier entre $(-\pi,~\pi)$}
    \label{figdes:2} 
\end{center}
\end{figure} 

Cómo se puede ver en la Figura \ref{figdes:2}, gracias a la calculadora DESMOS se puede ver un poco cómo es el comportamiento de la función, con ello, se comienza a calcular la serie de Fourier para $F(t)$, por tanto sus coeficientes son:

Para $a_0$ teniendo $\omega_f = 1$:
\begin{align*}
    a_0 &= \frac{\cancel{2}}{\cancel{2}\pi} \int_{-\frac{\cancel{2}\pi}{\cancel{2}}}^{\frac{\cancel{2}\pi}{\cancel{2}}} F(t) \,dt\\
    &= \frac{1}{\pi} \int_{-\pi}^{\pi} e^{t} \,dt\\
    &= \frac{1}{\pi} \left[e^{t} |_{-\pi}^{\pi}\right]\\
    &= \frac{1}{\pi} \left[e^{\pi} - e^{-\pi}\right]\\
    &= \frac{1}{\pi} \left(\frac{e^{2\pi}-1}{e^{\pi}}\right)
\end{align*}
Resultando
\begin{equation}
    \textcolor{rojo}{a_0 = \frac{e^{2\pi} - 1}{\pi e^{\pi}}}
    \label{res:22}
\end{equation}

Para $a_n$ teniendo $\omega_f = 1$:
\begin{align*}
    a_n &= \frac{\cancel{2}}{\cancel{2}\pi} \int_{-\frac{\cancel{2}\pi}{\cancel{2}}}^{\frac{\cancel{2}\pi}{\cancel{2}}} F(t) \cos\left(nt\right) \,dt\\
    &= \frac{1}{\pi} \int_{-\pi}^{\pi} e^{t} \cos\left(nt\right) \,dt
\end{align*}
Sustituyendo al usar integración por partes
\begin{align*}
    u&=\cos\left(nt\right)           &  dv&=e^{t} dt\\
    du&=-n\sin\left(nt\right)dt         &  v&=e^{t}
\end{align*}
Entonces
\begin{align*}
    a_n = \frac{1}{\pi} \left[e^{t} \cos\left(nt\right) + n\int e^{t} \sin\left(nt\right) \,dt\right]_{-\pi}^{\pi}
\end{align*}
Sustituyendo nuevamente al usar integración por partes
\begin{align*}
    u&=\sin\left(nt\right)           &  dv&=e^{t} dt\\
    du&=n\cos\left(nt\right)dt         &  v&=e^{t}
\end{align*}
Entonces
\begin{align*}
    a_n &= \frac{1}{\pi} \left[e^{t} \cos\left(nt\right)\right]_{-\pi}^{\pi} +n\\
    &\hphantom{=}~~ \left[\left(e^{t} \sin\left(nt\right) - n\int e^{t} \cos\left(nt\right) \,dt\right)\right]_{-\pi}^{\pi}\\
    &= \frac{1}{\pi} \left[e^{t} \cos\left(nt\right) + ne^{t}\sin\left(nt\right)\right]_{-\pi}^{\pi} -\\
    &\hphantom{=}~~ n^{2}\int_{-\pi}^{\pi} e^{t} \cos\left(nt\right) \,dt\\
    &= \frac{1}{\pi} \left[e^{t} \cos\left(nt\right) + ne^{t} \sin\left(nt\right)\right]_{-\pi}^{\pi} \cdot \frac{1}{n^{2}+1}\\
    &\textup{Tras desarrollar los límites, teniendo que }\\
    &\sin\left(n\left(\pm \pi\right)\right) = 0~~\mathrm{y}~\cos\left(n\left(\pm \pi\right)\right)\\
    &= \frac{1}{\pi} \left[e^{\pi}\left(-1\right)^{n}-e^{-\pi}\left(-1\right)^{n}\right] \cdot \frac{1}{n^{2}+1}\\
    &= \frac{e^{2\pi}-1}{\pi e^{\pi}} \left(\frac{\left(-1\right)^{n}}{n^{2}+1}\right)
\end{align*}
Resultando
\begin{equation}
    \textcolor{rojo}{a_n = \frac{e^{2\pi}-1}{\pi e^{\pi}} \left(\frac{\left(-1\right)^{n}}{n^{2}+1}\right)}
    \label{res:23}
\end{equation}

Para $b_n$ teniendo $\omega_f = 1$:
\begin{align*}
    b_n &= \frac{\cancel{2}}{\cancel{2}\pi} \int_{-\frac{\cancel{2}\pi}{\cancel{2}}}^{\frac{\cancel{2}\pi}{\cancel{2}}} F(t) \sin\left(nt\right) \,dt\\
    &= \frac{1}{\pi} \int_{-\pi}^{\pi} e^{t} \sin\left(nt\right) \,dt
\end{align*}
Sustituyendo al usar integración por partes
\begin{align*}
    u&=\sin\left(nt\right)           &  dv&=e^{t} dt\\
    du&=n\cos\left(nt\right)dt         &  v&=e^{t}
\end{align*}
Entonces
\begin{align*}
    b_n = \frac{1}{\pi} \left[e^{t} \sin\left(nt\right) - n\int e^{t} \cos\left(nt\right) \,dt\right]_{-\pi}^{\pi}
\end{align*}
Sustituyendo nuevamente al usar integración por partes
\begin{align*}
    u&=\cos\left(nt\right)           &  dv&=e^{t} dt\\
    du&=-n\sin\left(nt\right)dt         &  v&=e^{t}
\end{align*}
Entonces
\begin{align*}
    b_n &= \frac{1}{\pi} \left[e^{t} \sin\left(nt\right)\right]_{-\pi}^{\pi} -n\\
    &\hphantom{=}~~ \left[\left(e^{t} \cos\left(nt\right) + n\int e^{t} \sin\left(nt\right) \,dt\right)\right]_{-\pi}^{\pi}\\
    &= \frac{1}{\pi} \left[e^{t} \sin\left(nt\right) - ne^{t}\cos\left(nt\right)\right]_{-\pi}^{\pi} -\\
    &\hphantom{=}~~ n^{2}\int_{-\pi}^{\pi} e^{t} \sin\left(nt\right) \,dt\\
    &= \frac{1}{\pi} \left[e^{t} \sin\left(nt\right) - ne^{t} \cos\left(nt\right)\right]_{-\pi}^{\pi} \cdot \frac{1}{n^{2}+1}\\
    &\textup{Tras desarrollar los límites, teniendo que }\\
    &\sin\left(n\left(\pm \pi\right)\right) = 0~~\mathrm{y}~\cos\left(n\left(\pm \pi\right)\right)\\
    &= \frac{1}{\pi} \left[-ne^{\pi}\left(-1\right)^{n}+ne^{-\pi}\left(-1\right)^{n}\right] \cdot \frac{1}{n^{2}+1}\\
    &= \frac{e^{2\pi}-1}{\pi e^{\pi}} \left(\frac{-n\left(-1\right)^{n}}{n^{2}+1}\right)
\end{align*}
Resultando
\begin{equation}
    \textcolor{rojo}{b_n = \frac{e^{2\pi}-1}{\pi e^{\pi}} \left(\frac{-n\left(-1\right)^{n}}{n^{2}+1}\right)}
    \label{res:24}
\end{equation}
Usando (\ref{res:22}), (\ref{res:23}) y (\ref{res:24}), reemplazando en la forma general de la serie trigonométrica de Fourier y factorizando, resulta la serie de Fourier para $F(t)$:
\begin{equation*}
    \textcolor{rojo}{F(t) = \frac{e^{2\pi}-1}{\pi e^{\pi}} \left[\frac{1}{2} + \sum_{n = 1}^{\infty} \frac{\left(-1\right)^{n}}{n^{2}+1} \left(\cos\left(nt\right) - n\sin\left(nt\right)\right) \right]}
\end{equation*}

$c.~F(t) = F(t+2\pi) = \ln|t|;~(-\pi,~\pi)$

Usando el software \textit{DESMOS}, la gráfica para $F(t)$ junto a su serie de Fourier correspondiente evaluada hasta $n=15$ es:

\begin{figure}[H]
\begin{center}
    \includegraphics[scale=0.40]{desmos-graph3.png}
    \caption{$F(t) = \ln |t|$ con su serie de Fourier entre $(-\pi,~\pi)$}
    \label{figdes:3}
\end{center}
\end{figure} 
    
Cómo se puede ver en la Figura \ref{figdes:3}, gracias a la calculadora DESMOS se puede ver cómo es el comportamiento de la función, teniendo una asíntota en 0 y cada $2\pi$ tal como lo menciona su periodo. Debido a esto, al realizar las respectivas integraciones para la función, su complejidad es alta llegando a tener que utilizar la función $\Gamma$ o integrales especiales para seno y coseno, por ende, solamente se va a enseñar el comando usado en el software y se va a asumir la serie de Fourier partiendo de la Figura \ref{figdes:3}.

Teniendo $T_f = 2\pi$ y $\omega_f = 1$:
\begin{figure}[H]
\begin{center}
    \includegraphics[scale=0.60]{capture.png}
    \caption{Código usado para hallar la serie de Fourier de $F(t)=|t|$}
    \label{capture}
\end{center}
\end{figure} 


\begin{enumerate}[leftmargin=15pt, resume]
    \item La serie de Fourier de la función $F(t)$ definida por:
\end{enumerate}

\begin{equation}
    F(t) =
    \begin{cases}
    -1,&\text{si } \sfrac{-\pi}{2} \leq t < 0\\
    \hphantom{-}1,&\text{si } \hspace{18pt} 0 \leq t < \sfrac{\pi}{2}
    \end{cases}
\label{enu:3}
\end{equation}
\hspace{12pt} y
\begin{equation*}
    F(t) = F\left(t + \pi\right)
\end{equation*}
\hspace{12pt} Viene dada por:
\begin{equation}
    F(t) = \sum_{\mathrm{impar}}^{\infty} \frac{4}{n\pi} \sin\left(2nt\right)
    \label{enu:4}
\end{equation}
\hspace{12pt} A partir del resultado anterior, demostrar que:
\begin{equation}
    \frac{\pi}{4} = 1 - \frac{1}{3} + \frac{1}{5} - \frac{1}{7} + \frac{1}{9} - \hdots
    \label{enu:5}
\end{equation}

\subsubsection*{\textbf{Solución:}}

Para la solución de este ejercicio nos ayudamos de un software para graficar como lo es \textit{DESMOS}, con ello poder analizar la función ya graficada y conocer su forma, la gráfica en cuestión es:

\begin{figure}[H]
\begin{center}
\includegraphics[scale=0.40]{desmos-graphpr3.png}
\caption{Gráfica de la función mostrada en (\ref{enu:3})}
\label{fig:graf31}
\end{center}
\end{figure}

Analizando la simetría de la gráfica vista en Figura \ref{fig:graf31}, se puede deducir que la función es \textit{impar}, lo que implica que nos podemos ahorrar algunos cálculos realizando la serie de Fourier para una función impar.

Teniendo en que si $F(t)$ es impar, sus coeficientes de la serie de Fourier son:
\begin{equation}
    a_0 = a_n = 0
    \label{eq:5}
\end{equation}
y
\begin{equation}
    b_n = \frac{4}{T_f} \int_{0}^{\sfrac{T_f}{2}} F(t)\sin\left(n \omega_f t\right) \,dt
    \label{eq:6} 
\end{equation}
Teniendo que $T_f = \pi$, $\omega_f = 2$ y usando (\ref{eq:5}) y (\ref{eq:6}) para hallar los coeficientes, resulta:

Para $a_0$ y $a_n$, según (\ref{eq:5}), $a_0 = a_n = 0$.

Para $b_n$:
\begin{align*}
    b_n &= \frac{4}{\pi} \int_{0}^{\sfrac{\pi}{2}} 1 \sin\left(2nt\right) \,dt\\
    &= \frac{4}{\pi} \int_{0}^{\sfrac{\pi}{2}} \sin\left(2nt\right) \,dt\\
    &= \frac{4}{\pi} \left[-\frac{1}{2n} \left(\cos\left(2nt\right) |^{\sfrac{\pi}{2}}_{0}\right)\right]\\
    &= \frac{4}{\pi} \left[-\frac{1}{2n} \left(\cos\left(\cancel{2}n\frac{\pi}{\cancel{2}}\right) - \cos\left(2n(0)\right)\right)\right]\\
    &\mathrm{Ya~que}~\cos\left(n\pi\right) = \left(-1\right)^{n} \mathrm{~y~} \cos\left(0\right) = 1\\
    &= \frac{4}{\pi} \left[-\frac{1}{2n}\left(\left(-1\right)^{n} - 1\right)\right]
\end{align*}
Al igual que con el con el caso del itemize 1. de esta misma sección, solo interesa que $n$ sea impar, ya que se realizó el análisis solo se mencionará que $n$ es impar en la serie de Fourier final.

Resultando
\begin{equation}
    b_n = \frac{4}{n\pi},~\mathrm{donde}~n~\mathrm{es~impar}
    \label{res:7}
\end{equation}
De esta manera, la serie de Fourier para (\ref{enu:3}) es:
\begin{equation}
    \textcolor{rojo}{F(t) = \frac{4}{\pi} \sum_{\mathrm{impar}}^{\infty} \frac{\sin\left(2nt\right)}{n}}
    \label{res:8}
\end{equation}
Donde se puede ver la clara similitud a (\ref{enu:4}), demostrando así que esa es su forma.

Para demostrar (\ref{enu:5}), se decide evaluar la serie de Fourier mostrada en (\ref{res:8}) en $t = \sfrac{\pi}{4}$. La razón es por la forma de la demostración que se está requiriendo, ya que se tiene $\sfrac{4}{\pi}$ como factor constante en la serie, se debe evaluar la función dentro del intervalo $0 \leq t < \sfrac{\pi}{2}$ para que esta sea igual a $1$, se usa $\sfrac{\pi}{4}$ estando este dentro del intervalo y satisfaciendo la demostración requerida. Por tanto
\begin{align*}
    F\left(\frac{\pi}{4}\right) &= \frac{4}{\pi} \sum_{impar}^{\infty} \frac{\sin\left(\cancel{2}n\left(\frac{\pi}{\cancel{4}}\right)\right)}{n}\\
    1 &= \frac{4}{\pi} \sum_{impar}^{\infty} \frac{\sin\left(n\frac{\pi}{2}\right)}{n}\\
    \frac{\pi}{4} &= \sum_{impar}^{\infty} \frac{\sin\left(n\frac{\pi}{2}\right)}{n}
\end{align*}
Ya que $\sin\left(n\frac{\pi}{2}\right) = \left(-1\right)^{n}$
\begin{equation}
    \frac{\pi}{4} = \sum_{n = 1}^{\infty} \frac{\left(-1\right)^{n}}{2n-1}
    \label{eq:7}
\end{equation}
Como se puede ver en (\ref{eq:7}), se diferenció a $n$, debido a que esta ya no es necesariamente impar, sino puede ser cualquier valor siendo $n > 0,~n\in \mathbb{Z}$, por lo tanto, solamente se reemplaza la $n$ del denominador en su forma impar ya que esta no fue afectada por la función.

Evaluando $n$ en la serie, resultaría:
\begin{equation}
    \textcolor{rojo}{\frac{\pi}{4} = 1 - \frac{1}{3} + \frac{1}{5} - \frac{1}{7} + \frac{1}{9} - \hdots} + \frac{(-1)^{n+1}}{2(n+1) - 1}
    \label{eq:8}
\end{equation}
Donde se puede ver la similitud de los 5 primeros términos de la serie entre (\ref{eq:8}) y (\ref{enu:5}), demostrando y con ello solucionando el ejercicio.

\begin{enumerate}[leftmargin=15pt, resume]
    \item Encontrar las componentes par e impar de las siguientes funciones:
    \begin{enumerate}
        \item $F(t) = e^{t}$
        \item $F(t) = \frac{t+1}{t-1}$
    \end{enumerate}
\end{enumerate}

\subsubsection*{\textbf{Solución:}}
Teniendo en cuenta que una función $G(x)$ puede ser expresada como una combinación de sus componentes par e impar, de la forma:
\begin{equation}
    G(x) = G_p(x) + G_i(x)
    \label{eq:9}
\end{equation}
Las ecuaciones que se emplearán para hallar cada componente, son:
\begin{equation}
    G_p(x) = \frac{1}{2}\left(G(x) + G(-x)\right)
    \label{eq:10}
\end{equation}
\begin{equation}
    G_i(x) = \frac{1}{2}\left(G(x) - G(-x)\right)
    \label{eq:11}
\end{equation}

$a.~ F(t) = e^{t}$

Usando (\ref{eq:10}) y (\ref{eq:11}) para hallar las componentes de $F(t)$, resulta:

Para $F_p(t)$:
\begin{align*}
    F_p(t) &= \frac{1}{2}\left(e^{t} + e^{-t}\right)\\
    &= \frac{e^{t} + e^{-t}}{2}
\end{align*}
Recordando la forma del coseno hiperbólico ($\cosh$), la cual es:
\begin{equation}
    \cosh = \frac{e^{t} + e^{-t}}{2}
    \label{eq:12}
\end{equation}
Reemplazando (\ref{eq:12}), resulta
\begin{equation}
    \textcolor{rojo}{F_p(t) = \cosh}
    \label{res:9}
\end{equation}
Para $F_i(t)$:
\begin{align*}
    F_i(t) &= \frac{1}{2}\left(e^{t} - e^{-t}\right)\\
    &= \frac{e^{t} - e^{-t}}{2}
\end{align*}
Recordando la forma del seno hiperbólico ($\sinh$), la cual es:
\begin{equation}
    \sinh = \frac{e^{t} - e^{-t}}{2}
    \label{eq:13}
\end{equation}
Reemplazando (\ref{eq:13}), resulta
\begin{equation}
    \textcolor{rojo}{F_i(t) = \sinh}
    \label{res:10}
\end{equation}
Para finalizar, usando (\ref{eq:9}), la función $F(t)$ expresada en sus componentes par e impar es:
\begin{equation}
    \textcolor{rojo}{e^{t} = \cosh + \sinh}
    \label{res:11}
\end{equation}

$b.~F(t) = \frac{t+1}{t-1}$

Usando (\ref{eq:10}) y (\ref{eq:11}) para hallar las componentes de $F(t)$, resulta:

Para $F_p(t)$:
\begin{align*}
    F_p(t) &= \frac{1}{2} \left(\frac{t+1}{t-1} + \frac{-t+1}{-t-1}\right)\\
    &= \frac{1}{2} \left(\frac{t+1}{t-1} + \frac{t-1}{t+1}\right)\\
    &= \frac{1}{2} \left(\frac{(t+1)^{2} + (t-1)^{2}}{t^{2} - 1}\right)\\
    &= \frac{1}{2} \left(\frac{t^{2}~ \cancel{+~2t} + 1 + t^{2}~ \cancel{-~2t} + 1}{t^{2} - 1}\right)\\
    &= \frac{1}{2} \left(\frac{2t^{2} + 2}{t^{2} - 1}\right)\\
    &= \frac{1}{\cancel{2}} \left(\frac{\cancel{2}\left(t^{2} + 1\right)}{t^{2} - 1}\right)\\
    &= \frac{t^{2} + 1}{t^2 - 1}
\end{align*}
Resultando
\begin{equation}
    \textcolor{rojo}{F_p(t) = \frac{t^{2} + 1}{t^2 - 1}}
    \label{res:12}
\end{equation}
Para $F_i(t)$:
\begin{align*}
    F_i(t) &= \frac{1}{2} \left(\frac{t+1}{t-1} - \frac{-t+1}{-t-1}\right)\\
    &= \frac{1}{2} \left(\frac{t+1}{t-1} - \frac{t-1}{t+1}\right)\\
    &= \frac{1}{2} \left(\frac{(t+1)^{2} - (t-1)^{2}}{t^{2} - 1}\right)\\
    &= \frac{1}{2} \left(\frac{\cancel{t^{2}} + 2t~ \cancel{+~1}~ \cancel{-~t^{2}} + 2t~ \cancel{-~1}}{t^{2} - 1}\right)\\
    &= \frac{1}{\cancel{2}} \left(\frac{\cancel{4}t}{t^{2} - 1}\right)\\
    &= \frac{2t}{t^{2} - 1}
\end{align*}
Resultando
\begin{equation}
    \textcolor{rojo}{F_i(t) = \frac{2t}{t^{2} - 1}}
    \label{res:13}
\end{equation}
Para finalizar, usando (\ref{eq:9}), la función $F(t)$ expresada en sus componentes par e impar es:
\begin{equation}
    \textcolor{rojo}{F(t) = \frac{t^{2} + 1}{t^2 - 1} + \frac{2t}{t^{2} - 1}}
    \label{res:14}
\end{equation}
Para demostrarlo solo es cuestión de realizar la suma, desarrollar unos productos notables y factorizar, por tanto:
\begin{align*}
    F(t) &= \frac{t^{2} + 1}{t^{2} - 1} + \frac{2t}{t^{2} - 1}\\
    &= \frac{t^{2} + 2t + 1}{t^{2} - 1}\\
    &= \frac{\left(t + 1\right)^{\cancel{2}}}{\cancel{(t + 1)}(t - 1)}\\
    &= \frac{t + 1}{t - 1}
\end{align*}
Demostrando así que efectivamente (\ref{res:12}) y (\ref{res:13}) son las componentes par e impar de $F(t)$, respectivamente.

\begin{enumerate}[leftmargin=15pt, resume]
    \item Identifique si la función periódica definida por $F(t)=|t|$ para $(-\pi,~\pi)$ y $F(t)=F(t+2\pi)$, es par o impar. Dependiendo de su respuesta, calcule los coeficientes de Fourier que corresponda y encuentre su serie de Fourier.
\end{enumerate}

\begin{figure}[H]
\begin{center}
\includegraphics[scale=0.28]{Fig1.png}
\caption{Gráfica de $F(t)=|t|$ entre $(-\pi,~\pi)$ con periodo $T_f = 2\pi$}
\label{figenu:2}
\end{center}
\end{figure}

\subsubsection*{\textbf{Solución:}}

La función $F(t)=|t|$ entre $(-\pi,~\pi)$, también puede ser escrita como una función a trozos, de la forma
\begin{equation}
    F(t) =
    \begin{cases}
    -t,&\text{si } -\pi \leq t < 0\\
    \hphantom{-}t,&\text{si } \hphantom{-\pi} 0 \leq t \leq \pi
    \end{cases}
    \label{eq:14}
\end{equation}
Y analizando la forma de la función basándose en la Figura \ref{figenu:2}, se puede ver que la función es \textbf{par}, lo que nos ahorra un cálculo, el cuál sería $b_n$, esto gracias a las siguientes relaciones:

Si $F(t)$ es par:
\begin{equation}
    a_0 = \frac{4}{T_f} \int_{0}^{\sfrac{T_f}{2}} F(t) \,dt
    \label{eq:15}
\end{equation}
\begin{equation}
    a_n = \frac{4}{T_f} \int_{0}^{\sfrac{T_f}{2}} F(t)\cos\left(n \omega_f t\right) \,dt
    \label{eq:16}
\end{equation}
y $b_n = 0$.

Por lo tanto, se puede hacer la serie de Fourier teniendo en cuenta solamente los coeficientes $a_0$ y $a_n$, y el límite de integración necesario. Teniendo $T_f = 2\pi$ y $\omega_f = 1$ la serie de Fourier resultaría:

Para $a_0$:
\begin{align*}
    a_0 &= \frac{\cancel{4}}{\cancel{2}\pi} \int_{0}^{\frac{\cancel{2}\pi}{\cancel{2}}} F(t) \,dt\\
    &= \frac{2}{\pi} \int_{0}^{\pi} t \,dt\\
    &= \frac{2}{\pi} \left(\frac{t^{2}}{2} |_{0}^{\pi} \right)\\
    &= \frac{\cancel{2}}{\cancel{\pi}} \left(\frac{\pi^{\cancel{2}}}{\cancel{2}}\right)\\
    &= \pi
\end{align*}
Resultando
\begin{equation}
    \textcolor{rojo}{a_0 = \pi}
    \label{res:15}
\end{equation}
Para $a_n$:
\begin{align*}
    a_n &= \frac{\cancel{4}}{\cancel{2}\pi} \int_{0}^{\frac{\cancel{2}\pi}{\cancel{2}}} F(t) \cos\left(nt\right) \,dt\\
    &= \frac{2}{\pi} \int_{0}^{\pi} t \cos\left(nt\right) \,dt
\end{align*}
Sustituyendo al usar integración por partes
\begin{align*}
    u&=t           &  dv&=\cos\left(nt\right) dt\\
    du&=dt         &  v&=\frac{1}{n} \sin\left(nt\right)
\end{align*}
Entonces
\begin{align*}
    a_n &= \frac{2}{\pi} \left[t\frac{1}{n} \sin\left(nt\right) - \int \frac{1}{n}\sin\left(nt\right) \,dt\right]_{0}^{\pi}\\
    &= \frac{2}{\pi} \left[\frac{t}{n} \sin\left(nt\right) - \frac{1}{n}\int \sin\left(nt\right) \,dt\right]_{0}^{\pi}\\
    &= \frac{2}{\pi} \left[\frac{t}{n} \sin\left(nt\right) - \frac{1}{n}\left(-\frac{1}{n} \cos\left(nt\right)\right)\right]_{0}^{\pi}\\
    &= \frac{2}{\pi} \left[\frac{t}{n} \sin\left(nt\right) + \frac{1}{n^{2}} \cos\left(nt\right)\right]_{0}^{\pi}\\
    &= \frac{2}{\pi} \left[\frac{nt \sin\left(nt\right) + \cos\left(nt\right)}{n^{2}}\right]_{0}^{\pi}\\
    &= \frac{2}{\pi} \cdot \left(\frac{n\pi \sin\left(n\pi\right) + \cos\left(n\pi\right)}{n^{2}}\right) -\\ 
    &\hphantom{=}~ \left(\frac{n(0) \sin\left(n(0)\right) + \cos\left(n(0)\right)}{n^{2}}\right)
\end{align*}
Ya que $\sin\left(n\pi\right) = \sin\left(0\right) = 0,~ \cos\left(n\pi\right) = \left(-1\right)^{n}$ y $\cos\left(0\right) = 1$,
\begin{align*}
    \hphantom{a_0} &= \frac{2}{\pi} \left(\frac{\left(-1\right)^n - 1}{n^{2}}\right) \textup{, tomando $n$ impar}\\
    &= \frac{2}{\pi} \left(\frac{-2}{\left(2n-1\right)^{2}}\right)\\
    &= -\frac{4}{\left(2n-1\right)^{2} \pi}
\end{align*}
Resultando
\begin{equation}
    \textcolor{rojo}{a_n = -\frac{4}{\left(2n-1\right)^{2} \pi}}
    \label{res:16}
\end{equation}
Reemplazando (\ref{res:15}) y (\ref{res:16}) en la ecuación general para la serie trigonométrica de Fourier:
\begin{equation}
    \textcolor{rojo}{F(t) = \frac{\pi}{2} - \frac{4}{\pi} \sum_{n = 1}^{\infty} \frac{\cos\left(\left(2n-1\right)t\right)}{\left(2n-1\right)^2}}
    \label{res:17.1}
\end{equation}
La cual también puede ser escrita de la forma:
\begin{equation}
    \textcolor{rojo}{F(t) = \frac{\pi}{2} - \frac{4}{\pi} \sum_{\mathrm{impar}}^{\infty} \frac{\cos\left(nt\right)}{n^2}}
    \label{res:17.2}
\end{equation}

\end{multicols}

\end{document}